\chapter{Lesson 5: Numbers, Time, Telephone Numbers and Prices}

\section{Numbers 0--10}

Japanese has two sets of numerals. The Sino-Japanese set (borrowed from Chinese) is
used for counting, time, money, phone numbers and most counters. A native set
(\ja{ひとつ}, \ja{ふたつ}, \ldots) appears later for counting objects.

\begin{center}
\begin{tabular}{clll}
\toprule
 & Kanji & Hiragana & Romaji \\
\midrule
0 & \ja{〇} / \ja{零} & \ja{ゼロ}, \ja{れい} & zero, rei \\
1 & \ja{一} & \ja{いち} & ichi \\
2 & \ja{二} & \ja{に} & ni \\
3 & \ja{三} & \ja{さん} & san \\
4 & \ja{四} & \ja{よん}, \ja{し} & yon, shi \\
5 & \ja{五} & \ja{ご} & go \\
6 & \ja{六} & \ja{ろく} & roku \\
7 & \ja{七} & \ja{なな}, \ja{しち} & nana, shichi \\
8 & \ja{八} & \ja{はち} & hachi \\
9 & \ja{九} & \ja{きゅう}, \ja{く} & kyū, ku \\
10 & \ja{十} & \ja{じゅう} & jū \\
\bottomrule
\end{tabular}
\end{center}

\begin{remark}
4, 7 and 9 have two readings. \ja{し} (4) sounds like \ja{死} ``death'' and \ja{く} (9)
like \ja{苦} ``suffering'', so \ja{よん}, \ja{なな}, \ja{きゅう} are preferred in
general counting. Which reading a particular counter takes must be memorised (see the
time tables below). Hospitals and hotels sometimes skip room numbers with 4.
\end{remark}

\section{Larger numbers}

Numbers are built multiplicatively, exactly like reading the kanji:
\begin{itemize}
  \item 11 = \ja{じゅういち} (ten-one), 20 = \ja{にじゅう} (two-ten),
        35 = \ja{さんじゅうご} (three-ten-five).
  \item 100 = \ja{ひゃく}, 1{,}000 = \ja{せん}, 10{,}000 = \ja{まん} (\ja{いちまん}).
\end{itemize}

\begin{important}[Sound changes with hundreds and thousands]
\begin{center}
\begin{tabular}{llll}
\toprule
 & $\times 100$ & $\times 1000$ & \\
\midrule
1 & \ja{ひゃく} hyaku & \ja{せん} sen & (no \ja{いち}) \\
3 & \ja{さんびゃく} sanbyaku & \ja{さんぜん} sanzen & voicing \\
6 & \ja{ろっぴゃく} roppyaku & \ja{ろくせん} rokusen & \\
8 & \ja{はっぴゃく} happyaku & \ja{はっせん} hassen & \\
\bottomrule
\end{tabular}
\end{center}
All others are regular: \ja{よんひゃく}, \ja{ごひゃく}, \ja{ななひゃく},
\ja{きゅうひゃく}, \ja{にせん}, \ja{よんせん}, \ja{きゅうせん}.
\end{important}

\subsection{Counting by \ja{まん} (ten thousands)}
English groups digits in threes (thousand, million); Japanese groups them in fours
(\ja{まん} $=10^4$, \ja{おく} $=10^8$). So 20{,}000 is \ja{にまん} ``two \emph{man}'',
and 150{,}000 is \ja{じゅうごまん} ``fifteen \emph{man}''. It helps to rewrite the number
as 15{,}0000 before reading it.

\begin{center}
\begin{tikzpicture}[d/.style={draw, minimum width=1.1cm, minimum height=0.8cm, font=\large}]
  \foreach \x/\n in {0/1,1/2,2/3,3/4,4/5,5/6,6/7,7/8} \node[d] (c\x) at (\x*1.1,0) {\n};
  \draw[very thick, red!70!black] (3.85,-0.6) -- (3.85,0.6);
  \node[font=\small] at (3.85,0.9) {split every 4 digits};
  \node[font=\small, align=center, text width=4cm] at (1.65,-1.3) {\ja{せんにひゃくさんじゅうよんまん}\\1234 \emph{man}};
  \node[font=\small, align=center, text width=4cm] at (6.05,-1.3) {\ja{ごせんろっぴゃくななじゅうはち}\\5678};
\end{tikzpicture}
\end{center}

\section{Telephone numbers}

Phone numbers are read digit by digit, with \ja{の} \emph{no} for the dash.
Long vowels are often used for 2 and 5 to keep the rhythm: \ja{にい}, \ja{ごお}.
\begin{example}
519-884-xxxx (a Waterloo number):
\ja{ごいちきゅうの　はちはちよんの\ldots} \emph{go-ichi-kyū no hachi-hachi-yon no \ldots}
\end{example}

\begin{example}[Asking for a number]
\leavevmode
\begin{itemize}
  \item[A:] \ja{でんわばんごうはなんですか。}
        \emph{Denwa bangō wa nan desu ka.} What is your phone number?
  \item[B:] \ja{ろくよんななの、さんにいごおはちです。}
        \emph{Roku-yon-nana no, san-nii-gō-hachi desu.} It's 647-3258.
\end{itemize}
\end{example}

\section{Telling time}

\subsection{Hours: \ja{〜じ} (\ja{時})}
\begin{center}
\begin{tabular}{llll}
\toprule
Time & Reading & Time & Reading \\
\midrule
1:00 & \ja{いちじ} & 7:00 & \ja{しちじ} \textbf{(not \ja{なな})} \\
2:00 & \ja{にじ} & 8:00 & \ja{はちじ} \\
3:00 & \ja{さんじ} & 9:00 & \ja{くじ} \textbf{(not \ja{きゅう})} \\
4:00 & \ja{よじ} \textbf{(not \ja{よん})} & 10:00 & \ja{じゅうじ} \\
5:00 & \ja{ごじ} & 11:00 & \ja{じゅういちじ} \\
6:00 & \ja{ろくじ} & 12:00 & \ja{じゅうにじ} \\
\bottomrule
\end{tabular}
\end{center}
Note \ja{しちじ} uses \ja{しち} here; some speakers say \ja{ななじ} to avoid confusion
with \ja{いちじ}.

\subsection{Minutes: \ja{〜ふん}/\ja{〜ぷん} (\ja{分})}
The counter for minutes alternates between \emph{fun} and \emph{pun}, with small
\ja{っ} appearing:
\begin{center}
\begin{tabular}{llll}
\toprule
Minutes & Reading & Minutes & Reading \\
\midrule
1 & \ja{いっぷん} ippun & 6 & \ja{ろっぷん} roppun \\
2 & \ja{にふん} nifun & 7 & \ja{ななふん} nanafun \\
3 & \ja{さんぷん} sanpun & 8 & \ja{はっぷん} happun \\
4 & \ja{よんぷん} yonpun & 9 & \ja{きゅうふん} kyūfun \\
5 & \ja{ごふん} gofun & 10 & \ja{じゅっぷん} juppun \\
\bottomrule
\end{tabular}
\end{center}
Higher numbers follow the last digit: 15 \ja{じゅうごふん}, 20 \ja{にじゅっぷん},
45 \ja{よんじゅうごふん}. Half past is \ja{はん} \emph{han}.

\begin{center}
\begin{tikzpicture}
  \draw[thick] (0,0) circle (1.6);
  \foreach \h in {1,...,12} {
    \draw (90-\h*30:1.45) -- (90-\h*30:1.6);
    \node[font=\scriptsize] at (90-\h*30:1.25) {\h};
  }
  % 3:30 hour hand between 3 and 4, minute hand at 6
  \draw[very thick] (0,0) -- (90-3.5*30:0.8);
  \draw[thick] (0,0) -- (-90:1.35);
  \fill (0,0) circle (0.05);
  \node[below, font=\small] at (0,-1.75) {\ja{さんじはん} sanji han (3:30)};
  \begin{scope}[xshift=6.5cm]
  \draw[thick] (0,0) circle (1.6);
  \foreach \h in {1,...,12} {
    \draw (90-\h*30:1.45) -- (90-\h*30:1.6);
    \node[font=\scriptsize] at (90-\h*30:1.25) {\h};
  }
  % 9:10
  \draw[very thick] (0,0) -- (90-9.17*30:0.8);
  \draw[thick] (0,0) -- (90-2*30:1.35);
  \fill (0,0) circle (0.05);
  \node[below, font=\small] at (0,-1.75) {\ja{くじじゅっぷん} kuji juppun (9:10)};
  \end{scope}
\end{tikzpicture}
\end{center}

\subsection{Related time words}
\begin{center}
\begin{tabular}{lll}
\toprule
Japanese & Romaji & Meaning \\
\midrule
\ja{ごぜん} & gozen & a.m. (before the number: \ja{ごぜんくじ} 9 a.m.) \\
\ja{ごご} & gogo & p.m. (\ja{ごごさんじ} 3 p.m.) \\
\ja{いま} & ima & now \\
\ja{なんじ} & nanji & what time \\
\ja{〜ごろ} & goro & about (with times) \\
\ja{はん} & han & half (past) \\
\bottomrule
\end{tabular}
\end{center}

\begin{example}[What time is it?]
\leavevmode
\begin{itemize}
  \item[A:] \ja{すみません、いまなんじですか。} \emph{Sumimasen, ima nanji desu ka.}
        Excuse me, what time is it now?
  \item[B:] \ja{ごごにじはんです。} \emph{Gogo niji han desu.} It's 2:30 p.m.
\end{itemize}
\end{example}

\section{Prices and shopping numbers}

Prices use \ja{〜えん} (\ja{円}) ``yen''. Questions use \ja{いくら} ``how much''.
\begin{center}
\begin{tabular}{ll}
\toprule
Price & Reading \\
\midrule
¥100 & \ja{ひゃくえん} \\
¥300 & \ja{さんびゃくえん} \\
¥680 & \ja{ろっぴゃくはちじゅうえん} \\
¥1,000 & \ja{せんえん} \\
¥3,500 & \ja{さんぜんごひゃくえん} \\
¥8,000 & \ja{はっせんえん} \\
¥10,000 & \ja{いちまんえん} \\
¥24,800 & \ja{にまんよんせんはっぴゃくえん} \\
\bottomrule
\end{tabular}
\end{center}

\begin{example}
\ja{このかさはいくらですか。} \emph{Kono kasa wa ikura desu ka.} How much is this
umbrella?\\
\ja{せんにひゃくえんです。} \emph{Sen nihyaku en desu.} It's 1,200 yen.
\end{example}

\begin{aside}
A rough rule of thumb in 2017: ¥100 $\approx$ C\$1.15. Japan has a ¥1 coin (often
needed because of consumption tax) and the ¥10,000 bill is the largest common note.
Cash was still widely used then. When paying, put the money in the small tray
(\ja{トレー}) by the register rather than handing it directly to the clerk.
\end{aside}

\section{Age and years in school}
\begin{itemize}
  \item Age: \ja{〜さい} (\ja{歳}). Irregular: \ja{いっさい}, \ja{はっさい},
        \ja{じゅっさい}, and 20 is \ja{はたち}. ``How old?'' \ja{なんさい} or politely
        \ja{おいくつ}.
  \item Year in university: \ja{〜ねんせい}: \ja{いちねんせい} first-year student,
        \ja{にねんせい}, \ldots
\end{itemize}

\section*{Lesson 5 review}
\begin{itemize}
  \item Sino-Japanese numerals; 4, 7, 9 have two readings.
  \item Sound changes: \ja{さんびゃく}, \ja{ろっぴゃく}, \ja{はっぴゃく}; \ja{さんぜん},
        \ja{はっせん}.
  \item Big numbers are grouped by \ja{まん} (10{,}000).
  \item Phone numbers digit-by-digit with \ja{の}; time \ja{〜じ〜ふん/ぷん}, \ja{はん},
        \ja{ごぜん}/\ja{ごご}; 4 o'clock \ja{よじ}, 9 o'clock \ja{くじ}.
  \item Prices: \ja{〜えん}; ``how much?'' \ja{いくらですか}.
\end{itemize}
